% =====================================================================
%  Chapter 23 — The Proof in Natural Language, 7
% =====================================================================

\chapter{The Seesaw and the Averaging}
\label{ch:averaging2}

The argument must now average over the Hecke characters, and it cannot do so
directly: the period it has carries the character, and a period with a
character in it cannot be averaged. Two devices remove the obstruction in
sequence.

\section{The seesaw}Let $G$ be a unitary group over $E$ and let $G^0$ be the totally definite
unitary group obtained by restricting scalars to the maximal real subfield
$F$. For an Eisenstein series $E_1$ on $G$ and $E_2$ on $G^0$, the Rallis
seesaw formula is

\begin{quote}
\textbf{Proposition U (seesaw).} For suitable Eisenstein series,
\[
\int_{G^0\backslash G} \theta_A(g)\,E_1(g)\,dg
\;=\;
\frac{\hat g \cdot L_A\!\left(\tfrac12,\varepsilon\right)}{L(1,\varepsilon)}
\int_{G\backslash G} \theta_A(g)\,E_2(g)\,dg,
\]
where $\hat g$ is the number of components of $G^0\backslash G$ and $L_A$ is
the $L$-function of the form $A$.
\end{quote}

\begin{quote}
\textbf{Proposition V (why it helps).} The right-hand side is a period against
an Eisenstein series, and Eisenstein series are attached to Hecke characters.
The character has been traded for an $L$-value, and $L$-values average.
\end{quote}

This is the classical trade: the left-hand side has the character inside the
integral, where it is awkward; the right-hand side has it inside $L_A$, where
it is an $L$-value and can be summed over.

\section{The Bessel period}Apply the seesaw to the theta lift. A theta lift to a group of higher rank has
a period that is a product of ordinary Bessel functions: the Weil
representation restricted to a torus is one-dimensional, its $L$-function is a
product of Bessel $L$-functions, and the theta series is a sum over that
character.

\begin{quote}
\textbf{Proposition W (Bessel period).} For the lift of $\pi$ from
$\mathrm{U}(E,n)$ to $\mathrm{U}(E,d)$, the period has the form
\[
B_A(\varphi, F) \;=\; \sum_{\gamma} \varphi(\gamma)\,\theta_A(\gamma)\,F(\gamma),
\]
with $\varphi$ an automorphic form and $F$ a cuspidal form, and this
quantity does not depend on the Hecke character except through a fixed
character already extracted from the lift.
\end{quote}

The formal counterpart is \code{bB} in
\code{FunctionModelRallisFromInputs.lean}, with the lift operator \code{vth}
alongside it, and the non-vanishing condition \code{PerNZ}.

The character is extracted by the parity alternation: the character used at
stage $i$ is $\mu_s$ for even $i$ and $\mu_s\eta^{-1}$ for odd $i$, where
$\eta$ is the fixed character of the tower. The parity is forced --- the
weight of a theta series of rank $n$ is congruent to $n$ modulo $2$, so the
available weights alternate and the characters must alternate with them.

\section{The non-vanishing obligation}\begin{quote}
\textbf{Proposition X.} The period must not vanish: at some interior stage
$j < m$ the Bessel pairing is non-zero, or at the top stage $j = m$ the
top-level pairing is non-zero.
\end{quote}

This is \code{PerNZ}, and it is a real obligation rather than a formality: if
the period vanishes, the seesaw gives zero and there is nothing to average.

\section{The averaging}With the period independent of the character, the sum over the Hecke
characters is a sum over a fixed function. But averaging over a finite group
of a non-constant function loses: the average is smaller than the function's
value by a variance term, and that variance must be controlled.

The tool is Behrend's construction.

\begin{quote}
\textbf{Proposition Y (Behrend).} There is a set $T \subseteq \{1,\dots,n-1\}$
with no three distinct elements satisfying a non-trivial relation
$\alpha x_1 + \beta x_2 + \gamma x_3 = 0$ with $|\alpha|,|\beta|,|\gamma| \le K$,
and with
\[
|T| \;\ge\; n \exp\!\left(-C\sqrt{\log n \log(K+1)}\right)
\]
for an absolute constant $C$.
\end{quote}

The formal counterpart is \code{BehrendFree}, \code{BehrendBound} and
\code{ThickSet} in \code{Formalism/BehrendSets/Statement.lean}, with the
quantitative variance bounds \code{SumSq}, \code{SumLin}, \code{ShortOrdLarge}
and \code{ShortOrdSq}.

\section{The group, and why it is the right one}The group in the formalisation is
$(\mathbb{Z}/2^{s+2})^\times / \langle -1\rangle$, cyclic of order $2^s$,
and the relevant element is the Frobenius $\ell \mapsto \ell$-th power map,
whose order is what controls the growth of the corresponding $L$-function.

\begin{quote}
\textbf{Proposition Z.} There are absolute constants $40$ and $8$ such that
the sums of $\log \ell / \ell^2$ and of $\log \ell / \ell$ over a Behrend-free
set of $\ell$'s are bounded by $40\log M / M^2$ and $8(1+\log Y)^2/M$
respectively.
\end{quote}

This is the arithmetic content: the Behrend structure controls both how many
elements there are and how much each contributes.

\section{The equidistribution, which is the real payload}\begin{quote}
\textbf{Proposition AA (the equidistribution).} Let $T \subseteq
\{1,\dots,n-1\}$ be a Behrend set, and let $\mathrm{thick}(n,T,u)$ be its
translate by $u \in \mathbb{Z}/n$. Then for every finite set $I$ of
positions, each with a residue $\mathrm{pos}(i) \bmod n$,
\[
\sum_{u \in \mathbb{Z}/n}
\bigl|\{i \in I : \mathrm{pos}(i) \in \mathrm{thick}(n,T,u)\}\bigr|
\;=\; |I| \cdot |T|.
\]
\end{quote}

Each position lies in exactly $|T|$ of the $n$ translates. That is what makes
the average over the translates equal to the average over the group, with no
loss: the Behrend condition guarantees the translates overlap in a controlled
way, so the average is exact.

The formal counterpart is \code{TranslateMean}, with \code{OrbitOne} and
\code{WindowTwo} supplying the injectivity and the window bound that make the
identity hold.

\section{The remark's three pieces}The remark's items 11, 15 and 19 are now in place:

\begin{center}
\begin{tabular}{@{}>{\RaggedRight\arraybackslash\hyphenpenalty=10000\exhyphenpenalty=10000}p{0.32\textwidth}>{\RaggedRight\arraybackslash\hyphenpenalty=10000\exhyphenpenalty=10000}p{0.62\textwidth}@{}}
\toprule
\textbf{Item} & \textbf{What it supplies} \\
\midrule
11 & The group with small-order Frobenius elements, and the bounds on what
they contribute. \\
15 & The equidistribution of the translates --- the ``gcd'' and the
``extra averaging''. \\
19 & The precise choice of the Iwasawa tower, ``specifically using
Behrend''. \\
\bottomrule
\end{tabular}
\end{center}

And the remark's own footnote to item 19 --- ``This is the sort of trick that
intuitively to me is obviously going to be unnecessary'' --- says that the
author thought the Behrend construction would turn out not to be needed. The
development contains it anyway.

\begin{keytakeaways}
\item The seesaw trades a character-dependent period for a period against an
Eisenstein series, so that the character becomes an $L$-value and can be
summed over.
\item The theta lift's period is a product of Bessel functions and does not
carry the character, which the parity alternation with the fixed character
$\eta$ arranges.
\item The price is non-vanishing: the Bessel period must be non-zero, at an
interior stage or at the top.
\item Behrend's construction controls the variance of the average, and the
equidistribution \code{TranslateMean} --- each position lying in exactly
$|T|$ of the translates --- makes the averaging exact.
\end{keytakeaways}
